% eksperymentalne tłumaczenie na włoski

%

\subsection{Translacje zwane przesunięciami} % lang-pl
\index{translacja}% % lang-pl
\index{przesunięcie}% % lang-pl
\subsection{Traslazioni, dette anche spostamenti} % lang-it
\index{traslazione}% % lang-it
\index{spostamento}% % lang-it

\begin{definition}
[translacja] % lang-pl
[traslazione] % lang-it
    Niech $w = (w_x, w_y)$ będzie wektorem. % lang-pl
    Funkcję $T_w$ zadaną przez % lang-pl
    Sia $w = (w_x, w_y)$ un vettore. % lang-it
    La funzione $T_w$ definita da % lang-it
    \begin{equation}
        T_w(x, y) = (x + w_x, y + w_y)
    \end{equation}
    nazywamy translacją o wektor $w$. % lang-pl
    si chiama traslazione di vettore $w$. % lang-it
\end{definition}

Translacje to inaczej przesunięcia równoległe. % lang-pl
Rodzina wszystkich translacji tworzy grupę, zaś przesunięcie o wektor zerowy (odwzorowanie tożsamościowe) jest jej elementem neutralnym. % lang-pl
Le traslazioni sono anche dette spostamenti paralleli. % lang-it
L'insieme di tutte le traslazioni forma un gruppo e lo spostamento di vettore nullo (l'identità) ne è l'elemento neutro. % lang-it
Bartłomiej Bzdęga omówi w $\Delta_{23}^{9}$ podstawy rachunku wektorowego (wektory jako odcinki zorientowane i jako pary liczb, regułę trójkąta dla dodawania wektorów), a w $\Delta_{23}^{10}$ doda do tego iloczyn skalarny i jego zastosowanie do dowodzenia prostopadłości; oba artykuły zakończy zadaniami. % lang-pl
Bartłomiej Bzdęga tratterà in $\Delta_{23}^{9}$ i fondamenti del calcolo vettoriale (vettori come segmenti orientati e come coppie di numeri, regola del triangolo per la somma di vettori), e in $\Delta_{23}^{10}$ vi aggiungerà il prodotto scalare e la sua applicazione alla dimostrazione della perpendicolarità; entrambi gli articoli si concluderanno con degli esercizi. % lang-it

%